\section*{الفصل \ref{ch:b1:p-block-13-15} --- الفئة p (1): زمر البور والكربون والنيتروجين}

\begin{solution}{exo:b1:p-block-13-15:1}
\ce{NH4NO3}: \babelsublr{$-$III} في \ce{NH4+}، و\babelsublr{$+$V} في \ce{NO3-}. \ce{N2O5}: \babelsublr{$+$V}.
\ce{HNO2}: \babelsublr{$+$III}. \ce{N2H4}: \babelsublr{$-$II}. \ce{Mg3N2}: \babelsublr{$-$III}.
\end{solution}

\begin{solution}{exo:b1:p-block-13-15:2}
\ce{BF3}: البور بثلاث روابط وستة إلكترونات، مثلثي مستوٍ.
\ce{NH3}: ثلاث روابط \omterm{def:b1:lewis-resonance-vsepr:lewis}{وزوج حر}. يكوّن \omterm{def:b1:lewis-resonance-vsepr:lewis}{الزوج الحر} للنيتروجين رابطة تساندية
مع B: \ce{BF3} \omterm{def:b1:lewis-resonance-vsepr:lewis-acid}{حمض لويس}، و\ce{NH3} \omterm{def:b1:lewis-resonance-vsepr:lewis-acid}{قاعدة لويس}؛ وفي
المركب الإضافي تكون الذرتان كلتاهما رباعيتي الأوجه.
\end{solution}

\begin{solution}{exo:b1:p-block-13-15:3}
\ce{CO2 + 2NaOH -> Na2CO3 + H2O}; \ce{SiO2 + 2NaOH -> Na2SiO3 + H2O};
\ce{P4O10 + 12NaOH -> 4Na3PO4 + 6H2O}; \ce{Al2O3 + 6HCl -> 2AlCl3 + 3H2O};
\ce{Al2O3 + 2NaOH + 3H2O -> 2NaAl(OH)4}.
\end{solution}

\begin{solution}{exo:b1:p-block-13-15:4}
$\Delta_r G^\circ = 2 \times (-16.45) = \qty{-32.9}{kJ/mol}$، $\log K =
32.9/5.708 = 5.76$، $K = 10^{5.76}$.
\end{solution}

\begin{solution}{exo:b1:p-block-13-15:5}
حلقة من ست وحدات \ce{SiO3^2-}: \ce{Si6O18^12-}. السلسلة المزدوجة، لكل أربع
ذرات سيليكون: اثنتان برأسين متقاسمين (3 O، والشحنة $-2$ لكل منهما)، واثنتان
بثلاثة رؤوس ($2.5$ O، والشحنة $-1$ لكل منهما): \ce{Si4O11^6-}.
\end{solution}

\begin{solution}{exo:b1:p-block-13-15:6}
\qty{550}{kg} من \ce{Al(OH)3} تساوي $550/78.0 = \qty{7.05}{kmol}$، فتعطي
\qty{3.53}{kmol} من \ce{Al2O3}، أي \qty{360}{kg}؛ وبنسبة \qty{92}{\percent}:
\qty{331}{kg}.
\end{solution}

\begin{solution}{exo:b1:p-block-13-15:7}
للبور مدار فارغ يستقبل \omterm{def:b1:lewis-resonance-vsepr:lewis}{الزوج الحر} لأيون هيدروكسيد
مأخوذ من الماء: \ce{B(OH)3 + 2H2O <=> B(OH)4- + H3O+}. والبروتون
المحرَّر يأتي من جزيء ماء، لا من حمض البوريك.
\end{solution}

\begin{solution}{exo:b1:p-block-13-15:8}
ينتقل N من \babelsublr{$-$III} إلى \babelsublr{$+$II}، أي خمسة إلكترونات لكل \ce{NH3}؛ ويأخذ كل \ce{O2}
أربعة. $4 \times 5 = 5 \times 4$: \ce{4NH3 + 5O2 -> 4NO + 6H2O}.
\end{solution}

\begin{solution}{exo:b1:p-block-13-15:9}
\ce{NH3 + HNO3 -> NH4NO3}. الطن الواحد يساوي $10^6/80.0 = \qty{1.25e4}{mol}$؛
\qty{1.25e4}{mol} من الأمونيا للمعادلة و\qty{1.25e4}{mol} لصنع
حمض النتريك: $2.50 \times 10^4 \times 17.0 = \qty{425}{kg}$.
\end{solution}

\begin{solution}{exo:b1:p-block-13-15:10}
ذرات \omterm{def:b1:periodicity:table}{الدورة} الثانية صغيرة: فمداراتها \textit{p} تتداخل جيدًا جنبًا
إلى جنب وتكوّن روابط $\pi$ قوية، فيكون \ce{N#N} و\ce{O=C=O} جزيئين صغيرين
ثابتين. أما ذرات \omterm{def:b1:periodicity:table}{الدورة} الثالثة فأكبر؛ وتداخل $\pi$ فيها
ضعيف، فتفضّل روابط أحادية أكثر: رباعيات أوجه \ce{P4}،
وشبكة \ce{SiO4}.
\end{solution}

\begin{solution}{exo:b1:p-block-13-15:11}
\babelsublr{$E^\circ(\ce{PbO2}/\ce{Pb^2+}) = 1.46 > E^\circ(\ce{Cl2}/\ce{Cl-}) = 1.36$~V}:
\ce{PbO2 + 4H+ + 2Cl- -> Pb^2+ + Cl2 + 2H2O}. فالرصاص(IV) غير ثابت
مقارنةً بالرصاص(II) (الزوج الخامل)؛ أما السيليكون(II) فليس حالة ثابتة
و\ce{SiO2} ليس مؤكسدًا.
\end{solution}

\begin{solution}{exo:b1:p-block-13-15:12}
$160 \times 17.0/14.0 = \qty{194}{Mt}$ من الأمونيا، تحتوي $194 \times
3.0/17.0 = \qty{34}{Mt}$ من الهيدروجين، المصنوع أساسًا من الغاز الطبيعي
(الميثان وبخار الماء).
\end{solution}

\begin{solution}{pb:b1:p-block-13-15:1}
\textbf{1.} \babelsublr{:N\,$\equiv$\,N:} --- رابطة ثلاثية، قوية جدًا، بلا قطبية.
\textbf{2.} 0، \babelsublr{$-$III}، \babelsublr{$+$II}، \babelsublr{$+$IV}، \babelsublr{$+$V}، و\babelsublr{$-$III/$+$V} في \ce{NH4NO3}.
\textbf{3.} البكتيريا على جذور البقوليات، والبرق، وطريقة
هابر--بوش.
\textbf{4.} كسر الرابطة الثلاثية يحتاج إلى \omterm{def:b1:elementary-steps:catalyst}{حفّاز} لا تملكه النباتات؛
أما البكتيريا ذات إنزيم النيتروجيناز فتملكه.
\textbf{5.} النيتروجين يُختزل (من 0 إلى \babelsublr{$-$III})، والهيدروجين يتأكسد (من 0 إلى \babelsublr{$+$I}).
\textbf{6.} \ce{N2 + 3H2 <=> 2NH3}.
\textbf{7.} $K = 10^{2 \times 16.45/5.708} = 10^{5.76}$.
\textbf{8.} عند \qty{25}{\celsius} يكون التفاعل بطيئًا جدًا؛ \omterm{def:b1:elementary-steps:catalyst}{والحفّاز}
لا يعمل إلا ساخنًا، والظروف حل وسط يناقشه
مجلد السنة الثانية.
\textbf{9.} \qty{58.8}{kmol} من الأمونيا: \qty{29.41}{kmol} من \ce{N2}،
أي \qty{824}{kg}؛ و\qty{88.2}{kmol} من \ce{H2}، أي \qty{176}{kg}.
\textbf{10.} $V(\ce{N2}) = 29.4 \times 10^3 \times 8.314 \times 298.15/10^5 =
\qty{729}{m^3}$؛ والهواء $729/0.78 = \qty{935}{m^3}$.
\textbf{11.} من الغاز الطبيعي (إصلاح الميثان بالبخار)، مع
ناتج جانبي هو \ce{CO2}.
\textbf{12.} إذا حُقنت في التربة كانت سمادًا رخيصًا، لكنها غاز
سام؛ والأجسام الصلبة (اليوريا، ونترات الأمونيوم) أسهل تخزينًا ونثرًا.
\textbf{13.} N: \babelsublr{$-$III $\to$ $+$II} (5 إلكترونات)، O: \babelsublr{0 $\to$ $-$II} (4 لكل
\ce{O2}): \ce{4NH3 + 5O2 -> 4NO + 6H2O}.
\textbf{14.} \ce{2NO + O2 -> 2NO2}.
\textbf{15.} \ce{3NO2 + H2O -> 2HNO3 + NO}: تفاعل \omterm{def:b1:e-ph-diagrams:disproportionation}{عدم تناسب} (N من \babelsublr{$+$IV} إلى
\babelsublr{$+$V} و\babelsublr{$+$II}).
\textbf{16.} \ce{NH3 + 2O2 -> HNO3 + H2O}.
\textbf{17.} $2 \times 58.8 \times 32.0 = \qty{3.76}{t}$.
\textbf{18.} يجعل الأكسدة سريعة وانتقائية لأحادي أكسيد النيتروجين NO، بدل
\ce{N2} الأكثر ثباتًا.
\textbf{19.} \ce{NH3 + HNO3 -> NH4NO3}.
\textbf{20.} النصف.
\textbf{21.} \qty{29.4}{kmol} من \ce{NH4NO3}، أي \qty{2.35}{t}.
\textbf{22.} $28.0/80.0 = \qty{35}{\percent}$.
\textbf{23.} يحتوي مؤكسدًا (النترات) ومختزلًا (الأمونيوم) في
\omterm{def:b1:crystals-metals:crystal}{البلورة} نفسها: فإذا سُخّن أو خُلط بالمحروقات، أمكن أن يتفكك تفككًا انفجاريًا.
\textbf{24.} $58.8 \times 63.0 =$ \textbf{\qty{3.7}{t} من حمض النتريك لكل طن
من الأمونيا}.
\end{solution}
